% Generated draft: requires mathematical and source review.
\documentclass[11pt]{article}
\usepackage[margin=25mm]{geometry}
\usepackage{fontspec}
\setmainfont{Latin Modern Roman}
\usepackage{amsmath,amssymb,mathtools,unicode-math}
\setmathfont{Latin Modern Math}
\usepackage{enumitem,needspace,xurl,hyperref}
\setlength{\parindent}{0pt}
\setlength{\parskip}{5pt}
\setlength{\emergencystretch}{3em}
\newcommand{\N}{\mathbb{N}}
\newcommand{\Z}{\mathbb{Z}}
\newcommand{\Q}{\mathbb{Q}}
\newcommand{\R}{\mathbb{R}}
\newcommand{\C}{\mathbb{C}}
\newcommand{\F}{\mathbb{F}}
\newcommand{\A}{\mathbb{A}}
\newcommand{\PP}{\mathbb{P}}
\newcommand{\E}{\mathbb{E}}
\newcommand{\eps}{\varepsilon}
\newcommand{\dd}{\,\mathrm{d}}
\begin{document}
\section*{UnsolvedMath ID: 30003167}
\textbf{Source problem code:} catalogue entry
\subsection{Definitions and mathematical statement}
Let \(K\subseteq\C\) be a subfield, let \(X\) be a smooth projective variety over \(K\), and put \(X_{\C}=X\times_{\operatorname{Spec}K}\operatorname{Spec}\C\). Write \(2n+1\) for the odd cohomological degree in the source; the source does not explicitly specify the range of the integer \(n\). Set
\[
H_{\Z}=H^{2n+1}(X_{\C}(\C),\Z)/\operatorname{tors},\qquad
H_{\C}=H^{2n+1}(X_{\C}(\C),\C).
\]
Here \(\operatorname{tors}\) denotes the torsion subgroup, and \(H_{\Z}\) is identified with its image in \(H_{\C}\). The Hodge decomposition is
\[
H_{\C}=\bigoplus_{p+q=2n+1}H^{p,q}.
\]
The hypothesis that this cohomology has Hodge level \(1\) means, with the literal convention for level,
\[
\max\{\lvert p-q\rvert:H^{p,q}\ne0\}=1.
\]
Thus the only possible nonzero Hodge summands are \(H^{n+1,n}\) and \(H^{n,n+1}\). The Hodge-level condition is understood on the torsion-free cohomology; torsion carries no additional Hodge summands.

Let \(F^{\bullet}\) be the Hodge filtration. Define the intermediate Jacobian by
\[
J^{2n+1}(X_{\C})=
\frac{H_{\C}}{F^{n+1}H_{\C}+H_{\Z}},
\qquad
F^{n+1}H_{\C}=\bigoplus_{p\ge n+1}H^{p,2n+1-p}.
\]
Under the stated Hodge-level hypothesis and projectivity assumption, this complex torus is an abelian variety.

The question is whether, for every such \(K\), \(X\), and degree \(2n+1\), there exist an abelian variety \(A\) over \(K\) and an isomorphism of complex abelian varieties
\[
A\times_{\operatorname{Spec}K}\operatorname{Spec}\C
\;\cong\;J^{2n+1}(X_{\C}).
\]
This is the meaning of admitting a model over \(K\); the supplied question imposes no additional compatibility with a chosen polarization or Abel--Jacobi map.

The target is the whole intermediate Jacobian. The supplied literature assessment distinguishes it from the algebraic Abel--Jacobi image \(J_a^{2n+1}(X_{\C})\), namely the image of algebraically trivial cycles of codimension \(n+1\) under the Abel--Jacobi map. According to that supplied assessment, the cited distinguished-model result concerns this image and does not establish descent of the whole intermediate Jacobian in the stated generality. This assessment is supplied dated curation, not an independently verified literature conclusion.
\subsection{Short English statement}
For a smooth projective variety over a subfield \(K\subseteq\C\), assume its indicated odd-degree cohomology has Hodge level \(1\). Must the entire associated intermediate-Jacobian abelian variety be the complex base change of an abelian variety over \(K\)? The supplied record labels this question open and distinguishes it from descent of the algebraic Abel--Jacobi image.
\subsection{Sources}
\begin{itemize}
\item Supplied local UnsolvedMath registry record, permanent ID \(30003167\), title \emph{Descent of Intermediate Jacobians}, source problem code \texttt{catalogue entry}, in the Oberwolfach Reports Open Problems collection. No individual proposer is identified. The record gives proposed year \(2016\), status \texttt{open}, and statement status \texttt{corrected\_verified}; these are supplied registry labels. Its verification field reports manual semantic, typographic, and LaTeX review, which is not independently confirmed here.
\item Supplied report citation: \emph{Algebraic K-theory and Motivic Cohomology (2017)}, Oberwolfach Reports, DOI \texttt{10.4171/owr/2016/31}, pp.~1753--1807. \url{https://doi.org/10.4171/owr/2016/31}. The record describes this report as documenting workshop discussions and open problems; no exact internal page or named contributor is supplied for this question.
\item Achter--Casalaina-Martin--Vial, \emph{Distinguished models}, as labeled in the supplied literature sources. \url{https://arxiv.org/abs/1611.07471}.
\item Achter--Casalaina-Martin--Vial, \emph{Descending cohomology}, as labeled in the supplied literature sources. \url{https://arxiv.org/abs/1410.5376}.
\item Literature assessment dated exactly \texttt{2026-08-21}: the supplied curation labels the question open and states that the known descent theorem gives a distinguished model for the algebraic Abel--Jacobi image, rather than a model for the whole intermediate-Jacobian abelian variety. This is supplied dated curation, not an independently verified assessment. The source itself requests independent verification of time-sensitive status and cited claims.
\end{itemize}
\end{document}
